%
% eigenwerte.tex -- die drei Fälle sigma < 0, = 0, > 0
%
% basiert auf tikztemplate.tex
% (c) 2021 Prof Dr Andreas Müller, OST Ostschweizer Fachhochschule
%
\documentclass[tikz]{standalone}
\usepackage{amsmath}
\usepackage{times}
\usepackage{txfonts}
\usepackage{pgfplots}
\usepackage{csvsimple}
\usetikzlibrary{arrows,intersections,math,calc}
\definecolor{darkred}{rgb}{0.8,0,0}
\definecolor{darkgreen}{rgb}{0,0.6,0}
\begin{document}
    \def\skala{1}
    \begin{tikzpicture}[>=latex,thick,scale=\skala]

% sigma < 0 -- klingt ab
        \node[left] at (0.1, 2) {$\sigma < 0$:};
        \draw[->] (0.2,2) -- (3.4,2);
        \draw[->] (0.2,1.3) -- (0.2,2.7);
        \draw[domain=0:3, samples=100, color=blue]
        plot (\x+0.2, {2 + 0.5*exp(-1.2*\x)*sin(360*\x)});
        \draw[domain=0:3, samples=100, color=darkred, dashed]
        plot (\x+0.2, {2 + 0.5*exp(-1.2*\x)});
        \node[right] at (3.5, 2) {stabil};

% sigma = 0 -- Dauerschwingung
        \node[left] at (0.1, 0.5) {$\sigma = 0$:};
        \draw[->] (0.2,0.5) -- (3.4,0.5);
        \draw[->] (0.2,-0.1) -- (0.2,1.1);
        \draw[domain=0:3, samples=100, color=blue]
        plot (\x+0.2, {0.5 + 0.38*sin(360*\x)});
        \draw[color=darkred, dashed]
        (0.2,0.88) -- (3.4,0.88);
        \node[right] at (3.5, 0.5) {indifferent};

% sigma > 0 -- wächst
        \node[left] at (0.1, -1.0) {$\sigma > 0$:};
        \draw[->] (0.2,-1.0) -- (3.4,-1.0);
        \draw[->] (0.2,-1.6) -- (0.2,-0.3);
        \draw[domain=0:2.1, samples=100, color=blue]
        plot (\x+0.2, {-1 + 0.12*exp(0.9*\x)*sin(360*\x)});
        \draw[domain=0:2.1, samples=100, color=darkred, dashed]
        plot (\x+0.2, {-1 + 0.12*exp(0.9*\x)});
        \node[right] at (3.5, -1.0) {instabil};

    \end{tikzpicture}
\end{document}